\subsection{环的更换}

命题2. 设 A、B 为两个环， \(\rho: \mathbf{B} \to \mathbf{A}\) 为一个环同态，且 E（相应地，F）为右（相应地，左）A-模。则存在且仅存在一个 Z-线性映射

\[
\phi : \rho_ {*} (\mathbf {E}) \otimes_ {\mathbf {B}} \rho_ {*} (\mathbf {F}) \rightarrow \mathbf {E} \otimes_ {\mathbf {A}} \mathbf {F}\tag{6}
\]

使得对所有 \(x \in \mathbf{E}\) 与 \(y \in \mathbf{F}\) ， \(\phi\) 在元素 \(x \otimes y\) （它属于 \(\rho_{*}(\mathbf{E}) \otimes_{\mathbf{B}} \rho_{*}(\mathbf{F})\) ）处的像是元素 \(x \otimes y\) （它属于 \(\mathbf{E} \otimes_{\mathbf{A}} \mathbf{F}\) ） ；这个 \(\mathbf{Z}\) -线性映射是满射的。

我们考虑映射 \((x,y)\mapsto x\otimes y\) ，即从 \(\rho_{\star}(E)\times \rho_{*}(F)\) 到 \(E\otimes_A F\) 的映射；它是 Z-双线性的，并且对所有 \(\beta \in B\) ，根据定义有 \((x\rho (\beta))\otimes y = x\otimes (\rho (\beta)y)\) ，因此第1号的条件 (1) 成立，从而得到 \(\phi\) 的存在性与唯一性（第1号，命题1）。后一断言由元素 \(x\otimes y\) 生成 Z-模 \(E\otimes_A F\) 这一事实得出。

映射 (6) 称为典范映射。

推论。设 \(\mathfrak{a}\) 是A的一个双边理想，使得 \(\mathfrak{a}\) 包含于 E 的零化子与 F 的零化子中，从而 E（相应地 F）具有典范的左（相应地右） \((A/\mathfrak{a})\)-模结构（§ 1，第 12 号）。于是典范同态 (6)

\[
\phi : \mathbf {E} \otimes_ {\mathbf {A}} \mathbf {F} \rightarrow \mathbf {E} \otimes_ {\mathbf {A} / \mathfrak {a}} \mathbf {F}
\]

对应于典范同态 \(\rho: \mathbf{A} \to \mathbf{A} / \mathfrak{a}\) 是恒等映射。

对所有 \(\bar{\alpha} \in \mathbf{A} / \mathfrak{a}\) 、所有 \(x \in \mathbf{E}\) 以及所有 \(y \in \mathbf{F}\) ，有 \(x\bar{\alpha} = x\alpha\) （相应地， \(\bar{\alpha}y = \alpha y\) ），其中 \(\alpha\) 取遍满足 \(\rho(\alpha) = \bar{\alpha}\) 的一切元素。若 \(\mathbf{C} = \mathbf{Z}^{(\mathbb{E} \times \mathbb{F})}\) ，则 \(\mathbf{C}\) 中由元素 \((x\alpha, y) - (x, \alpha y)\) 生成的子模等于由元素 \((x\bar{\alpha}, y) - (x, \bar{\alpha}y)\) 生成的子模。

在命题2的假设和记号下，设 \(\mathbf{E}'\) 是一个右B-模， \(\mathbf{F}'\) 一个左B-模，并考虑两个半线性映射 \(u: \mathbf{E}' \to \mathbf{E}\) , \(v: \mathrm{F}' \to \mathrm{F}\) （它们相对于同态 \(\rho: \mathrm{B} \to \mathrm{A}\) ）； \(u\) （相应地 \(v\) ) 可被视为一个 B-线性映射 \(\mathrm{E}' \to \rho_*(\mathrm{E})\) （相应地 \(\mathrm{F}' \to \rho_*(\mathrm{F})\) )，从而得到一个 Z-线性映射 \(w: \mathrm{E}' \otimes_{\mathrm{B}} \mathrm{F}' \to \rho_*(\mathrm{E}) \otimes_{\mathrm{B}} \rho_*(\mathrm{F})\) ，使得 \(w(x' \otimes y') = u(x') \otimes v(y')\) ，其中 \(x' \in \mathrm{E}', y' \in \mathrm{F}'\) ；将典范映射 (6) 与这个映射复合，便得到一个 Z-线性映射 \(w': \mathrm{E}' \otimes_{\mathrm{B}} \mathrm{F}' \to \mathrm{E} \otimes_{\mathrm{A}} \mathrm{F}\) ，使得 \(w'(x' \otimes y') = u(x') \otimes v(y')\) ，其中 \(x' \in \mathrm{E}', y' \in \mathrm{F}'\) ；若不致引起混淆，这个映射通常就记作 \(u \otimes v\) 。显然， \((u, v) \mapsto u \otimes v\) 是一个 Z-双线性映射

\[
\operatorname{Hom} _ {\mathbf {B}} \left(\mathrm{E} ^ {\prime}, \rho_ {*} (\mathrm{E})\right) \times \operatorname{Hom} _ {\mathbf {B}} \left(\mathrm{F} ^ {\prime}, \rho_ {*} (\mathrm{F})\right)\rightarrow \operatorname{Hom} _ {\mathbf {Z}} \left(\mathrm{E} ^ {\prime} \otimes_ {\mathbf {B}} \mathrm{F} ^ {\prime}, \mathrm{E} \otimes_ {\mathbf {A}} \mathrm{F}\right).
\]

此外，若 C 是第三个环， \(\sigma: \mathrm{C} \to \mathrm{B}\) 是一个同态， \(\mathrm{E}''\) 是一个右C-模， \(\mathrm{F}''\) 是一个左C-模， \(u': \mathrm{E}'' \to \mathrm{E}'\) 和 \(v': \mathrm{F}'' \to \mathrm{F}'\) 是相对于 \(\sigma\) 的半线性映射，则

\[
(u \circ u ^ {\prime}) \otimes (v \circ v ^ {\prime}) = (u \otimes v) \circ (u ^ {\prime} \otimes v ^ {\prime}).
\]

